\section{Introduction}
\label{sec-intro}

Knowledge graphs such as DBpedia~\cite{auer2007dbpedia} and enterprise
graphs hold facts that users want to query in their own words.
Natural language interfaces translate a question into a structured
query that the graph can execute.
For relational data this has a long history, from interactive
interfaces that map parse trees onto SQL~\cite{li2014nalir} to recent
LLM-based systems~\cite{gao2024dailsql,li2024codes,ren2024purple,li2024dawn}.
For RDF, question answering has been cast as subgraph matching over the
graph~\cite{zou2014natural} and as template decomposition of the
question~\cite{zheng2018question}, and LLMs now generate SPARQL
directly~\cite{smeros2026sparqlllm,wang2026uniqgen,brei2025aruqula}.

Existing approaches leave three problems open.
\textbf{C1. Whole-query generation.}
Most LLM-based generators emit a complete query string.
A query that fails to parse, joins values of incompatible types, or
returns nothing is detected only after it is complete, and the repair
is another sample from the same model.
\textbf{C2. No static guarantees.}
A query can parse and still be wrong in ways that the schema already
rules out: a literal where the property expects an entity, a filter on
a variable that nothing binds, or a projected variable that is neither
grouped nor aggregated.
Such queries run and quietly return nothing, or an error, and the
generator learns little from either.
\textbf{C3. Search needs a reference.}
Constraint-guided search~\cite{wang2026uniqgen} refines a query by
adding and removing constraints with Chase and
Backchase~\cite{deutsch1999chase}, which gives soundness and
completeness guarantees, but relative to a reference answer set that
must be supplied or guessed, and over monotone conjunctive constraints,
with aggregation and ordering deferred.

\begin{example}
\label{exa-motivation}
Consider the question ``Which department does Ms.~Brant belong to?''
over the Text2SPARQL corporate graph~\cite{text2sparql}, asked in
German.
With Qwen3-14B~\cite{yang2025qwen3}, a one-shot generator wrote a query
over \kw{ecc{:}Employee} and \kw{foaf{:}name}~``Frau Brant'', terms the
graph does not use, so the query ran and returned nothing.
Constraint-guided search kept only the constraint
\kw{?e\ pv{:}memberOf\ ?d} and returned all six departments that have members.
Planning one operator at a time, the agent of
Section~\ref{sec-framework} first stated the goal: one answer of type
\kw{pv{:}Department}.
Looking up ``Frau Brant'' returned two employees, Karen and Sylvester
Brant; the agent bound Karen Brant, joined \kw{?e\ pv{:}memberOf\ ?d},
and saw one solution of the stated type: the correct department.
Each step was a single algebra operator, and each observation told the
agent what that operator did to the answer.
\end{example}

Database systems have handled a similar problem for decades.
A query optimizer does not guess a whole execution plan; it applies one
algebraic rewrite at a time and estimates its effect before keeping
it~\cite{selinger1979access}, and the quality of those estimates
largely decides the quality of the plan~\cite{leis2015good}.
The SPARQL algebra has a compositional semantics~\cite{perez2009semantics,sparql11},
so a query can be built and checked the same way.
This paper takes that view of natural language to SPARQL.

\stitle{Contributions}.
This paper introduces \hermes, a query-planning agent over a typed
formalization of the SPARQL algebra.
\begin{enumerate}
\item \textbf{A typed SPARQL algebra} (Section~\ref{sec-problem}).
  We treat query plans as terms of the SPARQL algebra and derive a
  typing context for their variables from the RDF
  schema~\cite{rdfs11}.
  Well-formedness and typing rules decide, without evaluating a plan,
  whether it is valid.
\item \textbf{Verified query planning} (Section~\ref{sec-problem}).
  We formulate natural language to SPARQL as planning: actions are
  algebra operators, observations report each operator's effect on the
  graph, and a plan is accepted only when it is well typed and
  satisfies a goal derived from the question.
\item \textbf{The \hermes system} (Section~\ref{sec-framework}).
  A verifier rejects ill-typed steps before they reach the graph store,
  and an LLM planner drives it.
  Existing generators are policies over the same action space.
\item \textbf{A reproducible evaluation} (Section~\ref{sec-exp}).
  We normalize six benchmarks and implement a scorer that reproduces
  all 5,024 published Text2SPARQL 2026 scores.
  With Qwen3 models of 4B to 14B parameters on the corporate benchmark,
  \hermes is the most accurate method at 14B, its accuracy grows
  fastest with model size, and the verifier cuts the items without a
  runnable query at every size.
\end{enumerate}
The code, data, configurations, and run outputs are available at the
artifact URL above.
